% =====================================================================
\chapter{Soluciones de los ejercicios}
\label{ap:soluciones}
% =====================================================================

Cada solución se presenta en cinco partes: datos, planteamiento, desarrollo, resultado y comprobación. En los pasos intermedios se redondea a cuatro cifras significativas salvo que se indique otra cosa; los resultados exactos se dan como fracción cuando existe. Con calculadora, $\log_2 t=\ln t/\ln 2=\ln t/0{,}6931$.

% ---------------------------------------------------------------------
\section{Tema 1}

\begin{solucion}{eje:1-aditiva}
\datos sucesos independientes $A$ y $B$ con $P(A)=p$ y $P(B)=q$, de modo que $P(A\cap B)=pq$.
\planteamiento la información de sucesos independientes debe ser aditiva mientras la probabilidad es multiplicativa; se busca la función que transforma productos en sumas (\cref{prop:caracterizacion}).
\pasos
\begin{enumerate}[label=\arabic*.]
  \item Si $A$ y $B$ son independientes, saber que ha ocurrido $B$ no cambia la sorpresa de $A$; la sorpresa conjunta debe ser la suma: $h(pq)=h(p)+h(q)$.
  \item El logaritmo cumple $\log(pq)=\log p+\log q$.
  \item Además, $h$ debe ser decreciente y valer~0 en $p=1$; $\log p$ es creciente y vale~0 en~1, de modo que hay que cambiarle el signo: $h(p)=-\log p$.
  \item Por la \cref{prop:caracterizacion}, toda solución decreciente es de la forma $-c\log_2p$ con $c\ge 0$; la base del logaritmo solo fija $c$, es decir, la unidad: base~2 da bits y base~$e$, nats.
\end{enumerate}
\resultado \textbf{$h(p)=-\log p$; la base solo elige la unidad de medida (bits o nats).}
\comprobacion dos monedas equilibradas: $h(1/4)=2$ bits $=h(1/2)+h(1/2)=1+1$.
\end{solucion}

\begin{solucion}{eje:1-archivos}
\datos dos archivos de $1$~MB $=8\cdot10^{6}$ bits: uno de ceros y otro de bits independientes y equiprobables.
\planteamiento el \cref{teo:cotas} (y el \cref{cor:bloques}) fija en $H$ bits por símbolo el mínimo de la longitud media de un código sin pérdidas.
\pasos
\begin{enumerate}[label=\arabic*.]
  \item Archivo de ceros: la fuente es degenerada, $H=0$ (\cref{prop:cotas-entropia}). Basta una descripción del tipo «$8\cdot10^{6}$ ceros», de unas decenas de bits.
  \item Archivo aleatorio: cada bit tiene distribución $(1/2,1/2)$ y $H=1$ bit por bit, el máximo posible para dos resultados.
  \item Ningún código sin pérdidas baja, en media, de $8\cdot10^{6}$ bits para el segundo archivo.
\end{enumerate}
\resultado \textbf{el archivo de ceros se comprime casi por completo ($H=0$); el aleatorio no admite compresión ($H=1$ bit por bit).}
\comprobacion en la práctica, un compresor de uso general aplicado a bits aleatorios produce un archivo algo mayor que el original, por sus cabeceras; es coherente con el \cref{eje:3-no-comprime}.
\end{solucion}

\begin{solucion}{eje:1-cambridge}
\datos (a) 256 valores equiprobables; (b) todos los seres humanos son mamíferos; (c) $p=(1/4,1/4,1/2)$; (d) la mediana divide la población en dos mitades, $p=(1/2,1/2)$.
\planteamiento $H=-\sum_xp(x)\log p(x)$; para la uniforme sobre $K$ valores, $H=\log K$.
\pasos
\begin{enumerate}[label=(\alph*)]
  \item $H=\log 256=\log 2^{8}=8$ bits.
  \item La distribución es degenerada: $P(\text{mamífero})=1$, luego $H=-1\cdot\log 1=0$ bits.
  \item $H=\tfrac14\cdot 2+\tfrac14\cdot 2+\tfrac12\cdot 1=0{,}5+0{,}5+0{,}5=1{,}5$ bits.
  \item $H=\tfrac12\cdot 1+\tfrac12\cdot 1=1$ bit.
\end{enumerate}
\resultado \textbf{(a) 8 bits; (b) 0 bits; (c) 1,5 bits; (d) 1 bit.}
\comprobacion (a) coincide con los 8 bits de un byte: un píxel uniforme no se puede comprimir. (c) es la distribución del \cref{ej:entropia-tres}.
\end{solucion}

\begin{solucion}{eje:1-binaria}
\datos $P(\text{cara})=1/4$; $H_b(\theta)=-\theta\log\theta-(1-\theta)\log(1-\theta)$ en $(0,1)$.
\planteamiento (a) sustituir; (b) derivar, usando $\frac{d}{d\theta}\log\theta=\frac{1}{\theta\ln 2}$, igualar a cero y estudiar la segunda derivada.
\pasos
\begin{enumerate}[label=\arabic*.]
  \item $-\log(1/4)=2$ y $-\log(3/4)=\log(4/3)=0{,}2877/0{,}6931=0{,}4150$.
  \item $H_b(1/4)=\tfrac14\cdot 2+\tfrac34\cdot 0{,}4150=0{,}5+0{,}3113=0{,}8113$ bits.
  \item Derivada: $\frac{d}{d\theta}\bigl(-\theta\log\theta\bigr)=-\log\theta-\frac{1}{\ln 2}$ y $\frac{d}{d\theta}\bigl(-(1-\theta)\log(1-\theta)\bigr)=\log(1-\theta)+\frac{1}{\ln 2}$.
  \item Sumando, $H_b'(\theta)=\log\frac{1-\theta}{\theta}$, que se anula si y solo si $\frac{1-\theta}{\theta}=1$, es decir, $\theta=1/2$.
  \item $H_b''(\theta)=\frac{1}{\ln 2}\Bigl(-\frac{1}{1-\theta}-\frac{1}{\theta}\Bigr)=-\frac{1}{\theta(1-\theta)\ln 2}<0$ en $(0,1)$: $H_b$ es estrictamente cóncava y el punto crítico es el único máximo, con $H_b(1/2)=1$ bit.
\end{enumerate}
\resultado \textbf{(a) $H_b(1/4)\approx 0{,}811$ bits; (b) $H_b'(\theta)=\log\frac{1-\theta}{\theta}$, con único máximo en $\theta=1/2$.}
\comprobacion $0<0{,}811<1$ y el valor coincide con el del \cref{ej:binaria}; además $H_b'>0$ para $\theta<1/2$ y $H_b'<0$ para $\theta>1/2$.
\end{solucion}

\begin{solucion}{eje:1-geometrica}
\datos $X$ uniforme en $\{1,\dots,8\}$; $P(Y=k)=2^{-k}$ para $k\ge 1$; $X$ e $Y$ independientes.
\planteamiento $H(X)=\log 8$; $H(Y)=\sum_k2^{-k}\cdot k$ porque $-\log2^{-k}=k$; por independencia, $p(x,y)=p(x)p(y)$ y $p(y\mid x)=p(y)$.
\pasos
\begin{enumerate}[label=\arabic*.]
  \item $H(X)=\log 8=3$ bits.
  \item $H(Y)=\sum_{k\ge 1}k\,2^{-k}$. Con la serie $\sum_{k\ge1}k\,t^{k}=t/(1-t)^{2}$ para $\abs{t}<1$ y $t=1/2$: $H(Y)=(1/2)/(1/4)=2$ bits.
  \item Por independencia, $H(Y\mid X=x)=H(Y)$ para cada $x$, luego $H(Y\mid X)=2$ bits.
  \item Regla de la cadena: $H(X,Y)=H(X)+H(Y\mid X)=3+2=5$ bits.
\end{enumerate}
\resultado \textbf{$H(X)=3$ bits $>H(Y)=2$ bits, aunque $Y$ tome infinitos valores; $H(X,Y)=5$ bits y $H(Y\mid X)=2$ bits.}
\comprobacion sumas parciales de $H(Y)$: $0{,}5$; $1$; $1{,}375$; $1{,}625$; $1{,}781$\ldots, que tienden a~2. El código prefijo \bin{0}, \bin{10}, \bin{110}, \ldots{} para $Y$ tiene longitud media exactamente~2 (\cref{teo:cota-inferior}, caso diádico).
\end{solucion}

\begin{solucion}{eje:1-funcion}
\datos $X$ con distribución $p$ y una función $g$; se escribe $S=\{x:p(x)>0\}$.
\planteamiento calcular $H(X,g(X))$ con la regla de la cadena (\cref{teo:cadena-entropia}) en los dos órdenes.
\pasos
\begin{enumerate}[label=\arabic*.]
  \item Primer orden: $H(X,g(X))=H(X)+H(g(X)\mid X)$. Conocido $X=x$, $g(X)$ vale $g(x)$ con probabilidad~1, luego $H(g(X)\mid X=x)=0$ para cada $x$ y $H(g(X)\mid X)=0$. Así, $H(X,g(X))=H(X)$.
  \item Segundo orden: $H(X,g(X))=H(g(X))+H(X\mid g(X))$.
  \item Igualando, $H(X)=H(g(X))+H(X\mid g(X))\ge H(g(X))$, porque $H(X\mid g(X))$ es una media de entropías, que son no negativas (\cref{prop:cotas-entropia}).
  \item Igualdad si y solo si $H(X\mid g(X))=0$, es decir, si para cada valor $y$ con probabilidad positiva la distribución de $X$ dado $g(X)=y$ es degenerada: hay un único $x\in S$ con $g(x)=y$. Esto es la inyectividad de $g$ sobre $S$.
  \item Interpretación: convertir el texto a minúsculas es una función no inyectiva («Casa» y «casa» van al mismo texto), de modo que la entropía del texto transformado es menor o igual que la del original; la información perdida es $H(X\mid g(X))$, la incertidumbre sobre las mayúsculas originales.
\end{enumerate}
\resultado \textbf{$H(g(X))\le H(X)$, con igualdad si y solo si $g$ es inyectiva sobre el soporte de $X$.}
\comprobacion si $X$ es uniforme en $\{1,2,3,4\}$ y $g(x)=x\bmod 2$, entonces $H(g(X))=1\le H(X)=2$, y $H(X\mid g(X))=1$ bit, que es la diferencia.
\end{solucion}

% ---------------------------------------------------------------------
\section{Tema 2}

\begin{solucion}{eje:2-vf}
\datos cuatro afirmaciones sobre la divergencia, la entropía cruzada, la entropía máxima y la información mutua.
\planteamiento contraejemplo numérico para (a); el \cref{teo:gibbs} y sus consecuencias para el resto.
\pasos
\begin{enumerate}[label=(\alph*)]
  \item En el \cref{ej:asimetria}, $\KL{p}{q}=0{,}322$ y $\KL{q}{p}=0{,}278$ bits. \textbf{Falsa.}
  \item $H(p,q)=H(p)+\KL{p}{q}$ y $\KL{p}{q}\ge0$ (\cref{prop:descomposicion}). \textbf{Verdadera.}
  \item Con $K=4$ símbolos, $H\le\log 4=2$ bits (\cref{prop:cotas-entropia}), y $3>2$. \textbf{Falsa.}
  \item $I(X;Y)=\KL{p(x,y)}{p(x)p(y)}$ se anula si y solo si $p(x,y)=p(x)p(y)$ para todo par (\cref{teo:mutua-positiva}). \textbf{Verdadera.}
\end{enumerate}
\resultado \textbf{(a) falsa; (b) verdadera; (c) falsa; (d) verdadera.}
\comprobacion para (c), 3 bits exigirían $2^{3}=8$ resultados equiprobables, el doble de los disponibles.
\end{solucion}

\begin{solucion}{eje:2-direcciones}
\datos la \cref{def:kl} y la \cref{obs:direcciones}.
\planteamiento (a) comparar con los axiomas de una distancia; (b) analizar qué términos de cada divergencia se hacen grandes.
\pasos
\begin{enumerate}[label=\arabic*.]
  \item (a) Una distancia es simétrica y cumple la desigualdad triangular. La entropía relativa no es simétrica (\cref{ej:asimetria}) ni cumple en general la desigualdad triangular; solo conserva la no negatividad y la anulación exacta en $p=q$ (\cref{teo:gibbs}). Por eso se llama divergencia.
  \item (b) Hacia delante, $\KL{p}{q}=\sum_xp(x)\log\frac{p(x)}{q(x)}$: los términos grandes aparecen donde $p(x)$ es apreciable y $q(x)$ es pequeña. Para evitarlos, $q$ debe repartir masa sobre las dos modas; con una sola moda, se sitúa entre ambas y es ancha: \emph{cubre} el soporte de $p$, aunque ponga masa en la zona intermedia, donde $p$ es pequeña.
  \item Hacia atrás, $\KL{q}{p}=\sum_xq(x)\log\frac{q(x)}{p(x)}$: los términos grandes aparecen donde $q(x)$ es apreciable y $p(x)$ es pequeña. A $q$ le conviene no poner masa en la zona intermedia y \emph{concentrarse} en una de las dos modas, ignorando la otra.
  \item El entrenamiento por máxima verosimilitud minimiza la entropía cruzada $H(p,q_\theta)$, que difiere de $\KL{p}{q_\theta}$ en la constante $H(p)$ (\cref{prop:descomposicion}): se minimiza la divergencia hacia delante.
\end{enumerate}
\resultado \textbf{(a) no es simétrica ni cumple la desigualdad triangular; (b) $\KL{p}{q}$ produce una $q$ que cubre ambas modas y $\KL{q}{p}$ una $q$ concentrada en una moda; la máxima verosimilitud minimiza $\KL{p}{q}$.}
\comprobacion la divergencia hacia atrás es la que aparece en la penalización del ajuste por preferencias, $\KL{\pi}{\pref}$, con la política que se optimiza en primer lugar (\cref{tema:5}).
\end{solucion}

\begin{solucion}{eje:2-kl-moneda}
\datos $p=(1/2,1/2)$, $q=(0{,}9;\ 0{,}1)$.
\planteamiento $H(p,q)=-\sum p\log q$; $\KL{p}{q}=H(p,q)-H(p)$; $\KL{q}{p}=H(q,p)-H(q)$.
\pasos
\begin{enumerate}[label=\arabic*.]
  \item $\log 0{,}9=-0{,}1520$ y $\log 0{,}1=-3{,}322$.
  \item $H(p,q)=\tfrac12\cdot 0{,}1520+\tfrac12\cdot 3{,}322=0{,}0760+1{,}661=1{,}737$ bits.
  \item $H(p)=1$ bit, luego $\KL{p}{q}=1{,}737-1=0{,}737$ bits.
  \item $H(q,p)=-0{,}9\log0{,}5-0{,}1\log 0{,}5=0{,}9+0{,}1=1$ bit.
  \item $H(q)=0{,}9\cdot 0{,}1520+0{,}1\cdot 3{,}322=0{,}1368+0{,}3322=0{,}4690$ bits.
  \item $\KL{q}{p}=1-0{,}4690=0{,}531$ bits.
\end{enumerate}
\resultado \textbf{$H(p,q)\approx 1{,}737$ bits; $\KL{p}{q}\approx 0{,}737$ bits; $\KL{q}{p}\approx 0{,}531$ bits.}
\comprobacion cálculo directo: $\KL{q}{p}=0{,}9\log 1{,}8+0{,}1\log 0{,}2=0{,}9\cdot 0{,}8480-0{,}1\cdot 2{,}322=0{,}7632-0{,}2322=0{,}531$.
\end{solucion}

\begin{solucion}{eje:2-conjunta}
\datos $p(0,0)=0{,}3$; $p(0,1)=0{,}2$; $p(1,0)=0{,}1$; $p(1,1)=0{,}4$.
\planteamiento marginales sumando filas y columnas; $I=H(X)+H(Y)-H(X,Y)$ (\cref{prop:mutua-identidades}); $H(Y\mid X)=H(X,Y)-H(X)$.
\pasos
\begin{enumerate}[label=\arabic*.]
  \item $p_X(0)=0{,}3+0{,}2=0{,}5$ y $p_X(1)=0{,}1+0{,}4=0{,}5$: $H(X)=1$ bit.
  \item $p_Y(0)=0{,}3+0{,}1=0{,}4$ y $p_Y(1)=0{,}2+0{,}4=0{,}6$.
  \item $H(Y)=0{,}4\cdot 1{,}322+0{,}6\cdot 0{,}7370=0{,}5288+0{,}4422=0{,}9710$ bits.
  \item $H(X,Y)=0{,}3\cdot 1{,}737+0{,}2\cdot 2{,}322+0{,}1\cdot 3{,}322+0{,}4\cdot 1{,}322=0{,}5211+0{,}4644+0{,}3322+0{,}5288=1{,}846$ bits.
  \item $I(X;Y)=1+0{,}9710-1{,}846=0{,}125$ bits.
  \item $H(Y\mid X)=1{,}846-1=0{,}846$ bits.
\end{enumerate}
\resultado \textbf{$H(X)=1$; $H(Y)\approx 0{,}971$; $H(X,Y)\approx 1{,}846$; $I(X;Y)\approx 0{,}125$; $H(Y\mid X)\approx 0{,}846$ bits.}
\comprobacion $H(Y)-H(Y\mid X)=0{,}971-0{,}846=0{,}125=I(X;Y)$, por la otra identidad de la \cref{prop:mutua-identidades}.
\end{solucion}

\begin{solucion}{eje:2-galletas}
\datos 6 galletas (3~C, 2~H, 1~A), extracción sin reemplazamiento; $\log 3=1{,}585$, $\log 5=2{,}322$.
\planteamiento construir la tabla conjunta $P(X=i,Y=j)=P(X=i)\,P(Y=j\mid X=i)$ y aplicar las definiciones de los temas~1 y~2.
\pasos
\begin{enumerate}[label=\arabic*.]
  \item $P(X)=(3/6,2/6,1/6)=(1/2,1/3,1/6)$.
  \item Condicionadas (quedan 5 galletas): $X=\mathrm{C}\Rightarrow Y\sim(2,2,1)/5$; $X=\mathrm{H}\Rightarrow(3,1,1)/5$; $X=\mathrm{A}\Rightarrow(3,2,0)/5$.
  \item Tabla conjunta en treintavos (filas $X=\mathrm{C},\mathrm{H},\mathrm{A}$; columnas $Y$ en el mismo orden):
  \[30P=\begin{pmatrix}6&6&3\\6&2&2\\3&2&0\end{pmatrix}.\]
  Las columnas suman $(15,10,5)$, luego $P(Y)=(1/2,1/3,1/6)$.
  \item $H(X)=H(Y)=\tfrac12\cdot 1+\tfrac13\log 3+\tfrac16\log 6=0{,}5+0{,}5283+0{,}4308=1{,}459$ bits.
  \item $H(X,Y)$: valores $1/5$ (tres veces), $1/10$ (dos veces) y $1/15$ (tres veces), con $-\log$ iguales a $2{,}322$, $3{,}322$ y $3{,}907$:
  $H(X,Y)=0{,}6\cdot 2{,}322+0{,}2\cdot 3{,}322+0{,}2\cdot 3{,}907=1{,}393+0{,}6644+0{,}7814=2{,}839$ bits.
  \item $H(Y\mid X)=H(X,Y)-H(X)=2{,}839-1{,}459=1{,}380$ bits, y por simetría de las marginales, $H(X\mid Y)=H(X,Y)-H(Y)=1{,}380$ bits.
  \item $I(X;Y)=H(Y)-H(Y\mid X)=1{,}459-1{,}380=0{,}079$ bits.
  \item (b) $P_{Y\mid X=\mathrm{H}}=(0{,}6;\ 0{,}2;\ 0{,}2)$ y $P_{Y\mid X=\mathrm{A}}=(0{,}6;\ 0{,}4;\ 0)$. El término de A en $\KL{P_{Y\mid X=\mathrm{H}}}{P_{Y\mid X=\mathrm{A}}}$ es $0{,}2\log(0{,}2/0)=+\infty$: la divergencia es infinita.
  \item $\KL{P_{Y\mid X=\mathrm{A}}}{P_{Y\mid X=\mathrm{H}}}=0{,}6\log 1+0{,}4\log 2+0=0{,}4$ bits.
  \item (c) Las ordenaciones distinguibles de los sabores son $6!/(3!\,2!\,1!)=720/12=60$, todas equiprobables: $H=\log 60=5{,}907$ bits.
\end{enumerate}
\resultado \textbf{$H(X)=H(Y)\approx 1{,}459$; $H(X,Y)\approx 2{,}839$; $H(Y\mid X)=H(X\mid Y)\approx 1{,}380$; $I(X;Y)\approx 0{,}079$ bits; $\KL{P_{Y\mid\mathrm{H}}}{P_{Y\mid\mathrm{A}}}=\infty$ y $\KL{P_{Y\mid\mathrm{A}}}{P_{Y\mid\mathrm{H}}}=0{,}4$ bits; el orden completo contiene $\log 60\approx 5{,}907$ bits.}
\comprobacion la tabla suma $30/30$. La divergencia infinita es la pérdida infinita del módulo~1: un modelo que da probabilidad~0 a algo que ocurre.
\end{solucion}

\begin{solucion}{eje:2-gibbs-ln}
\datos distribuciones $p$ y $q$; $S=\{x:p(x)>0\}$; $\ln t\le t-1$ para $t>0$, con igualdad solo en $t=1$; $u(x)=1/K$.
\planteamiento (a) aplicar la desigualdad a $t=q(x)/p(x)$ en $S$ y sumar con pesos $p(x)$; (b) desarrollar $\KL{p}{u}$.
\pasos
\begin{enumerate}[label=\arabic*.]
  \item (a) Si algún $x\in S$ tiene $q(x)=0$, la divergencia es $+\infty\ge 0$. En otro caso, los $x\notin S$ no contribuyen, y en nats $-\KL{p}{q}=\sum_{x\in S}p(x)\ln\frac{q(x)}{p(x)}$.
  \item Con $t=q(x)/p(x)$: $p(x)\ln\frac{q(x)}{p(x)}\le p(x)\Bigl(\frac{q(x)}{p(x)}-1\Bigr)=q(x)-p(x)$.
  \item Sumando sobre $S$: $-\KL{p}{q}\le\sum_{x\in S}q(x)-\sum_{x\in S}p(x)=\sum_{x\in S}q(x)-1\le 0$.
  \item Igualdad: el paso~2 exige $q(x)=p(x)$ en todo $S$ y el paso~3 exige $\sum_{x\in S}q(x)=1$, es decir, $q=0$ fuera de $S$. Juntas dan $q=p$. En bits basta dividir por $\ln 2>0$.
  \item (b) $\KL{p}{u}=\sum_xp(x)\log\frac{p(x)}{1/K}=\sum_xp(x)\log K+\sum_xp(x)\log p(x)=\log K-H(p)$.
  \item Por (a), $\log K-H(p)\ge 0$, con igualdad si y solo si $p=u$.
\end{enumerate}
\resultado \textbf{$\KL{p}{q}\ge 0$ con igualdad si y solo si $p=q$; en consecuencia, $H(p)\le\log K$ con igualdad solo para la uniforme.}
\comprobacion para $p=(1/2,1/4,1/4)$: $\KL{p}{u}=\log 3-1{,}5=1{,}585-1{,}5=0{,}085\ge 0$.
\end{solucion}

\begin{solucion}{eje:2-mutua}
\datos conjunta $p(x,y)$ con marginales $p(x)=\sum_yp(x,y)$ y $p(y)=\sum_xp(x,y)$.
\planteamiento desarrollar el logaritmo del cociente, agrupar por marginales y aplicar el \cref{teo:gibbs}.
\pasos
\begin{enumerate}[label=\arabic*.]
  \item Como $\log\frac{a}{bc}=\log a-\log b-\log c$,
  \[I(X;Y)=\sum_{x,y}p(x,y)\log p(x,y)-\sum_{x,y}p(x,y)\log p(x)-\sum_{x,y}p(x,y)\log p(y).\]
  \item El primer término es $-H(X,Y)$.
  \item En el segundo, se agrupa por $x$: $\sum_x\log p(x)\sum_yp(x,y)=\sum_xp(x)\log p(x)=-H(X)$; con el signo menos queda $+H(X)$.
  \item Análogamente, el tercer término es $+H(Y)$. Por tanto, $I(X;Y)=H(X)+H(Y)-H(X,Y)$.
  \item Por la regla de la cadena, $H(X,Y)=H(X)+H(Y\mid X)$; sustituyendo, $I(X;Y)=H(Y)-H(Y\mid X)$.
  \item $I(X;Y)$ es una entropía relativa, luego $I(X;Y)\ge 0$ y $H(Y\mid X)\le H(Y)$.
\end{enumerate}
\resultado \textbf{$I(X;Y)=H(X)+H(Y)-H(X,Y)=H(Y)-H(Y\mid X)\ge 0$, y condicionar no aumenta la entropía en media.}
\comprobacion con el \cref{eje:2-conjunta}: $H(Y\mid X)=0{,}846\le H(Y)=0{,}971$.
\end{solucion}

% ---------------------------------------------------------------------
\section{Tema 3}

\begin{solucion}{eje:3-dos-huffman}
\datos $p=(0{,}4;\ 0{,}2;\ 0{,}2;\ 0{,}1;\ 0{,}1)$ sobre $\{\mathrm{a},\mathrm{b},\mathrm{c},\mathrm{d},\mathrm{e}\}$ (\cref{ej:huffman}).
\planteamiento repetir el algoritmo cambiando el desempate de la tercera fusión y calcular las profundidades de las hojas.
\pasos
\begin{enumerate}[label=\arabic*.]
  \item Primeras fusiones, como en el texto: $\mathrm{d}+\mathrm{e}\to 0{,}2$; $\mathrm{b}+\mathrm{c}\to 0{,}4$. Quedan $\mathrm{a}$ ($0{,}4$), $\{\mathrm{b},\mathrm{c}\}$ ($0{,}4$) y $\{\mathrm{d},\mathrm{e}\}$ ($0{,}2$).
  \item El nodo menor es $\{\mathrm{d},\mathrm{e}\}$ y hay empate a $0{,}4$ para el segundo. En el texto se eligió $\mathrm{a}$; ahora se elige $\{\mathrm{b},\mathrm{c}\}$: $\{\mathrm{b},\mathrm{c}\}+\{\mathrm{d},\mathrm{e}\}\to 0{,}6$.
  \item Última fusión: $\mathrm{a}+0{,}6\to 1$. Profundidades: $\mathrm{a}$ a profundidad~1 y $\mathrm{b},\mathrm{c},\mathrm{d},\mathrm{e}$ a profundidad~3. Un código posible: $\mathrm{a}=\bin{0}$, $\mathrm{b}=\bin{100}$, $\mathrm{c}=\bin{101}$, $\mathrm{d}=\bin{110}$, $\mathrm{e}=\bin{111}$.
  \item Longitud media: $0{,}4\cdot 1+(0{,}2+0{,}2+0{,}1+0{,}1)\cdot 3=0{,}4+1{,}8=2{,}2$ bits.
\end{enumerate}
\resultado \textbf{sí: se obtienen las longitudes $\{2,2,2,3,3\}$ o $\{1,3,3,3,3\}$, ambas con longitud media $2{,}2$ bits; los dos códigos son óptimos (\cref{teo:huffman}).}
\comprobacion Kraft: $3\cdot2^{-2}+2\cdot2^{-3}=1$ y $2^{-1}+4\cdot2^{-3}=1$. Los códigos difieren en la varianza de la longitud: $0{,}16$ frente a $0{,}96$ bits$^{2}$.
\end{solucion}

\begin{solucion}{eje:3-diadica}
\datos $p=(1/2,1/4,1/8,1/8)$.
\planteamiento $H=\sum_ip_i(-\log p_i)$; código fijo de $\lceil\log 4\rceil$ bits; Huffman fusionando los dos nodos menores.
\pasos
\begin{enumerate}[label=\arabic*.]
  \item Sorpresas: $1$, $2$, $3$ y $3$ bits.
  \item $H=\tfrac12\cdot 1+\tfrac14\cdot 2+\tfrac18\cdot 3+\tfrac18\cdot 3=0{,}5+0{,}5+0{,}375+0{,}375=1{,}75$ bits.
  \item Código fijo: $\log 4=2$ bits por símbolo.
  \item Huffman: $\tfrac18+\tfrac18=\tfrac14$; quedan $(\tfrac12,\tfrac14,\tfrac14)$; $\tfrac14+\tfrac14=\tfrac12$; quedan $(\tfrac12,\tfrac12)$; fusión final.
  \item Palabras: \bin{0}, \bin{10}, \bin{110}, \bin{111}, de longitudes $1,2,3,3$; $L=1{,}75$ bits.
\end{enumerate}
\resultado \textbf{$H=1{,}75$ bits; código fijo de 2 bits; Huffman $\{\bin{0},\bin{10},\bin{110},\bin{111}\}$ con $L=1{,}75$ bits $=H$.}
\comprobacion Kraft: $2^{-1}+2^{-2}+2^{-3}+2^{-3}=1$. $L=H$ porque todas las probabilidades son potencias de $1/2$ (\cref{teo:cota-inferior}); es el código $C_4$ de la \cref{fig:arbol-kraft}.
\end{solucion}

\begin{solucion}{eje:3-huffman5}
\datos $p=(0{,}35;\ 0{,}25;\ 0{,}2;\ 0{,}15;\ 0{,}05)$; $H(p)=2{,}121$ bits.
\planteamiento fusionar los dos nodos menos probables; cada fusión suma un bit a las palabras de sus símbolos.
\pasos
\begin{enumerate}[label=\arabic*.]
  \item $0{,}05+0{,}15=0{,}2$. Quedan $0{,}35$; $0{,}25$; $0{,}2$; $0{,}2$.
  \item $0{,}2$ (símbolo) $+\,0{,}2$ (nodo) $=0{,}4$. Quedan $0{,}35$; $0{,}25$; $0{,}4$.
  \item $0{,}25+0{,}35=0{,}6$. Quedan $0{,}4$; $0{,}6$.
  \item $0{,}4+0{,}6=1$.
  \item Profundidades: $0{,}35\to 2$; $0{,}25\to 2$; $0{,}2\to 2$; $0{,}15\to 3$; $0{,}05\to 3$. Un código posible: \bin{00}, \bin{01}, \bin{10}, \bin{110}, \bin{111}.
  \item $L=0{,}7+0{,}5+0{,}4+0{,}45+0{,}15=2{,}2$ bits.
\end{enumerate}
\resultado \textbf{$L=2{,}2$ bits, $0{,}079$ por encima de $H$ y $0{,}8$ por debajo del código fijo de 3 bits.}
\comprobacion Kraft: $3\cdot 2^{-2}+2\cdot 2^{-3}=1$; y $2{,}121\le 2{,}2<3{,}121$, como exige el \cref{teo:cotas}.
\end{solucion}

\begin{solucion}{eje:3-unica}
\datos $C=\{\bin{00},\bin{11},\bin{0101},\bin{111},\bin{1010},\bin{100100},\bin{0110}\}$.
\planteamiento buscar prefijos, sumar $2^{-\ell}$ y buscar dos troceados distintos de una misma cadena.
\pasos
\begin{enumerate}[label=\arabic*.]
  \item \bin{11} es prefijo de \bin{111}: no es un código prefijo.
  \item Longitudes $2,2,4,3,4,6,4$: $\sum2^{-\ell}=\frac{16+16+4+8+4+1+4}{64}=\frac{53}{64}\approx 0{,}8281\le 1$. Cumple Kraft.
  \item La cadena \bin{11111} admite dos lecturas: $\bin{11}\,|\,\bin{111}$ y $\bin{111}\,|\,\bin{11}$.
\end{enumerate}
\resultado \textbf{no es prefijo; cumple Kraft ($53/64$); no es unívocamente descodificable: \bin{11111} es ambigua.}
\comprobacion Kraft garantiza que \emph{existe} un código prefijo con esas longitudes (\cref{teo:kraft}), no que este código concreto lo sea.
\end{solucion}

\begin{solucion}{eje:3-shannon}
\datos $p=(0{,}4;\ 0{,}2;\ 0{,}2;\ 0{,}1;\ 0{,}1)$; $H(p)=2{,}122$ bits; Huffman $L=2{,}2$ bits (\cref{ej:huffman}).
\planteamiento $\ell(x)=\lceil-\log p(x)\rceil$ (\cref{def:shannon}); las palabras se construyen como en la demostración del \cref{teo:kraft}.
\pasos
\begin{enumerate}[label=\arabic*.]
  \item $-\log 0{,}4=1{,}322\Rightarrow\ell=2$; $-\log 0{,}2=2{,}322\Rightarrow\ell=3$; $-\log 0{,}1=3{,}322\Rightarrow\ell=4$. Longitudes $(2,3,3,4,4)$.
  \item Kraft: $\frac14+\frac18+\frac18+\frac1{16}+\frac1{16}=0{,}625\le 1$.
  \item Palabras con $s_i=\sum_{j<i}2^{-\ell_j}$: $s=0,\ \tfrac14,\ \tfrac38,\ \tfrac12,\ \tfrac9{16}$, que dan \bin{00}, \bin{010}, \bin{011}, \bin{1000}, \bin{1001}.
  \item $L=0{,}4\cdot 2+0{,}2\cdot 3+0{,}2\cdot 3+0{,}1\cdot 4+0{,}1\cdot 4=0{,}8+1{,}2+0{,}8=2{,}8$ bits.
\end{enumerate}
\resultado \textbf{longitudes $(2,3,3,4,4)$, suma de Kraft $0{,}625$ y $L=2{,}8$ bits: cumple $2{,}122\le 2{,}8<3{,}122$, pero es $0{,}6$ bits peor que Huffman.}
\comprobacion la suma de Kraft $0{,}625<1$ indica que el código de Shannon desperdicia espacio: podrían acortarse palabras sin perder la propiedad de prefijo, que es lo que hace Huffman.
\end{solucion}

\begin{solucion}{eje:3-bloques}
\datos $p(\bin{a})=0{,}9$, $p(\bin{b})=0{,}1$; fuente sin memoria; $H_b(0{,}1)=0{,}469$ bits.
\planteamiento (a) dos símbolos; (b) las probabilidades de los pares son productos; Huffman sobre 4 símbolos; (c) \cref{cor:bloques} con $n=2$.
\pasos
\begin{enumerate}[label=\arabic*.]
  \item (a) Con dos símbolos, Huffman asigna \bin{0} y \bin{1}: 1 bit por símbolo.
  \item (b) Pares: $\bin{aa}=0{,}81$; $\bin{ab}=0{,}09$; $\bin{ba}=0{,}09$; $\bin{bb}=0{,}01$.
  \item Fusiones: $0{,}01+0{,}09=0{,}10$; $0{,}09+0{,}10=0{,}19$; $0{,}19+0{,}81=1$. Longitudes: $\bin{aa}\to 1$, un par de probabilidad $0{,}09\to 2$, el otro $\to 3$, $\bin{bb}\to 3$.
  \item $L_2=0{,}81\cdot 1+0{,}09\cdot 2+0{,}09\cdot 3+0{,}01\cdot 3=0{,}81+0{,}18+0{,}27+0{,}03=1{,}29$ bits por par, es decir, $0{,}645$ bits por símbolo.
  \item (c) La cota es $0{,}469\le L_2/2<0{,}469+0{,}5=0{,}969$, y $0{,}645$ la cumple.
\end{enumerate}
\resultado \textbf{(a) 1 bit por símbolo; (b) $0{,}645$ bits por símbolo; (c) se acerca a $H_b(0{,}1)=0{,}469$, dentro de la cota $[0{,}469;\ 0{,}969)$.}
\comprobacion Kraft para los pares: $2^{-1}+2^{-2}+2^{-3}+2^{-3}=1$. Para seguir acercándose a $0{,}469$ harían falta bloques mayores, con $2^{n}$ palabras: es la limitación que resuelve la codificación aritmética.
\end{solucion}

\begin{solucion}{eje:3-no-comprime}
\datos $C$ inyectiva de $\{\bin{0},\bin{1}\}^{+}$ en sí mismo; existe $s_0$ con $\abs{C(s_0)}<\abs{s_0}$.
\planteamiento por reducción al absurdo: suponer que ninguna cadena se alarga y contar cadenas de longitud acotada.
\pasos
\begin{enumerate}[label=\arabic*.]
  \item Supóngase que $\abs{C(s)}\le\abs{s}$ para toda $s$. Sea $n=\abs{s_0}$ y sea $A_m$ el conjunto de cadenas no vacías de longitud a lo sumo $m$; $\abs{A_m}=2+4+\dots+2^{m}=2^{m+1}-2$.
  \item Por la hipótesis, $C$ lleva $A_{n-1}$ en $A_{n-1}$. Como $C$ es inyectiva y $A_{n-1}$ es finito, $C$ restringida a $A_{n-1}$ es una biyección de $A_{n-1}$ en sí mismo.
  \item En particular, toda cadena de $A_{n-1}$ es imagen de alguna cadena de $A_{n-1}$.
  \item Pero $C(s_0)\in A_{n-1}$, porque $\abs{C(s_0)}<n$; por el paso~3, $C(s_0)=C(s)$ para alguna $s\in A_{n-1}$, y $s\ne s_0$ porque $\abs{s}<n=\abs{s_0}$.
  \item Esto contradice la inyectividad de $C$. Luego existe alguna cadena $s$ con $\abs{C(s)}>\abs{s}$.
\end{enumerate}
\resultado \textbf{todo compresor sin pérdidas que acorta alguna cadena alarga alguna otra.}
\comprobacion para $n=1$ el argumento dice que $C(s_0)$ tendría longitud~0, lo que no está permitido, y para $n=2$: las dos cadenas de longitud~1 deben ser imágenes de las dos cadenas de longitud~1, y no queda sitio para una tercera. Es la versión combinatoria del \cref{eje:1-archivos}.
\end{solucion}

% ---------------------------------------------------------------------
\section{Tema 4}

\begin{solucion}{eje:4-bpb-concepto}
\datos A: 2 bits por token y 2 bytes por token; B: 4 bits por token y 4 bytes por token.
\planteamiento separar lo que depende del tokenizador (el denominador «tokens») de lo que no depende (los bytes del texto y el tamaño del archivo comprimido).
\pasos
\begin{enumerate}[label=\arabic*.]
  \item La perplejidad por token es $2^{\text{bits por token}}$; depende de cuánto texto abarca cada token.
  \item Los bits totales $\sum_t-\log q(x_t\mid x_{1:t-1})$ son el tamaño del texto comprimido (\cref{teo:compresion}), una propiedad del texto completo.
  \item El número de bytes UTF-8 no depende del tokenizador; dividir el paso~2 por él da una tasa comparable, el bpb.
  \item A: perplejidad $2^{2}=4$ y $\bpb=2/2=1$. B: perplejidad $2^{4}=16$ y $\bpb=4/4=1$.
\end{enumerate}
\resultado \textbf{el bpb normaliza por los bytes del texto, que no dependen del tokenizador; A y B tienen perplejidades 4 y 16, pero los dos comprimen el texto a 1 bit por byte.}
\comprobacion para 1000 bytes, A usa 500 tokens $\times$ 2 bits y B usa 250 tokens $\times$ 4 bits: 1000 bits en ambos casos.
\end{solucion}

\begin{solucion}{eje:4-bac}
\datos $q=(1/2,1/4,1/4)$ para $(\mathrm{a},\mathrm{b},\mathrm{c})$; $F(\mathrm{a})=0$, $F(\mathrm{b})=1/2$, $F(\mathrm{c})=3/4$; mensaje «bac».
\planteamiento $a_t=a_{t-1}+w_{t-1}F(x_t)$, $w_t=w_{t-1}q(x_t)$ (\cref{def:intervalo}) con fracciones exactas; después buscar el intervalo diádico más largo contenido en el final.
\pasos
\begin{enumerate}[label=\arabic*.]
  \item Inicio: $a=0$, $w=1$.
  \item «b»: $a=0+1\cdot\tfrac12=\tfrac12$; $w=\tfrac14$. Intervalo $[1/2,3/4)$.
  \item «a»: $a=\tfrac12+\tfrac14\cdot 0=\tfrac12$; $w=\tfrac14\cdot\tfrac12=\tfrac18$. Intervalo $[1/2,5/8)$.
  \item «c»: $a=\tfrac12+\tfrac18\cdot\tfrac34=\tfrac{16}{32}+\tfrac{3}{32}=\tfrac{19}{32}$; $w=\tfrac18\cdot\tfrac14=\tfrac1{32}$. Intervalo $[19/32,20/32)=[0{,}59375;\ 0{,}625)$.
  \item Longitud ideal: $-\log(1/32)=5$ bits.
  \item Un intervalo diádico de 4 bits mide $1/16=2/32$ y no cabe en uno de anchura $1/32$. El intervalo final es él mismo diádico: $19/32=0{,}10011_2$, y $J(\bin{10011})=[19/32,20/32)$.
  \item El \cref{cor:longitud} garantiza $\lceil 5\rceil+1=6$ bits: $m=\lceil\tfrac{19}{32}\cdot 64\rceil=38$, palabra \bin{100110}.
\end{enumerate}
\resultado \textbf{intervalo $[19/32,\,5/8)$; longitud ideal 5 bits; palabra más corta \bin{10011} (5 bits); la construcción general garantiza 6 bits.}
\comprobacion $q(\text{bac})=\tfrac14\cdot\tfrac12\cdot\tfrac14=\tfrac1{32}$, que es la anchura final (\cref{prop:intervalos}).
\end{solucion}

\begin{solucion}{eje:4-descodificar}
\datos $q=(1/2,1/4,1/4)$ para $(\mathrm{a},\mathrm{b},\mathrm{c})$; $n=3$; palabra \bin{01110}.
\planteamiento $J(\bin{01110})=[14/32,15/32)=[0{,}4375;\ 0{,}46875)$; en cada paso se subdivide el intervalo actual y se elige el subintervalo que contiene a $J$.
\pasos
\begin{enumerate}[label=\arabic*.]
  \item Paso~1, sobre $[0,1)$: $\mathrm{a}=[0;\ 0{,}5)$, $\mathrm{b}=[0{,}5;\ 0{,}75)$, $\mathrm{c}=[0{,}75;\ 1)$. $J\subseteq[0;\ 0{,}5)$: primer símbolo \textbf{a}.
  \item Paso~2, sobre $[0;\ 0{,}5)$ con anchura $0{,}5$: $\mathrm{a}=[0;\ 0{,}25)$, $\mathrm{b}=[0{,}25;\ 0{,}375)$, $\mathrm{c}=[0{,}375;\ 0{,}5)$. $J\subseteq[0{,}375;\ 0{,}5)$: segundo símbolo \textbf{c}.
  \item Paso~3, sobre $[0{,}375;\ 0{,}5)$ con anchura $0{,}125$: $\mathrm{a}=[0{,}375;\ 0{,}4375)$, $\mathrm{b}=[0{,}4375;\ 0{,}46875)$, $\mathrm{c}=[0{,}46875;\ 0{,}5)$. $J=[0{,}4375;\ 0{,}46875)$: tercer símbolo \textbf{b}.
  \item $q(\text{acb})=\tfrac12\cdot\tfrac14\cdot\tfrac14=\tfrac1{32}$ y $-\log\tfrac1{32}=5$ bits.
\end{enumerate}
\resultado \textbf{el mensaje es «acb»; la palabra recibida tiene 5 bits, exactamente $-\log q(\text{acb})$.}
\comprobacion codificando «acb»: $a=0$, $w=1\to$ «a»: $[0,\tfrac12)$ $\to$ «c»: $a=\tfrac12\cdot\tfrac34=\tfrac38$, $w=\tfrac18$ $\to$ «b»: $a=\tfrac38+\tfrac18\cdot\tfrac12=\tfrac7{16}=\tfrac{14}{32}$, $w=\tfrac1{32}$. El intervalo es $[14/32,15/32)=J(\bin{01110})$.
\end{solucion}

\begin{solucion}{eje:4-abbb}
\datos $p(\mathrm{A})=1/4$, $p(\mathrm{B})=3/4$; $F(\mathrm{A})=0$, $F(\mathrm{B})=1/4$; mensaje «ABBB».
\planteamiento Huffman con dos símbolos gasta 1 bit por símbolo; aritmética con la \cref{def:intervalo}; longitud ideal $-\log p(\text{ABBB})$.
\pasos
\begin{enumerate}[label=\arabic*.]
  \item Huffman: $\mathrm{A}=\bin{0}$, $\mathrm{B}=\bin{1}$; «ABBB» $=\bin{0111}$, 4 bits.
  \item «A»: $[0,1/4)$, $w=1/4$.
  \item «B»: $a=0+\tfrac14\cdot\tfrac14=\tfrac1{16}$; $w=\tfrac14\cdot\tfrac34=\tfrac3{16}$. Intervalo $[1/16,1/4)$.
  \item «B»: $a=\tfrac1{16}+\tfrac3{16}\cdot\tfrac14=\tfrac{7}{64}$; $w=\tfrac{9}{64}$. Intervalo $[7/64,1/4)$.
  \item «B»: $a=\tfrac7{64}+\tfrac9{64}\cdot\tfrac14=\tfrac{37}{256}$; $w=\tfrac{27}{256}$. Intervalo $[37/256,64/256)=[0{,}1445;\ 0{,}25)$.
  \item Longitud ideal: $-\log\tfrac{27}{256}=8-\log 27=8-4{,}755=3{,}245$ bits.
  \item Intervalos diádicos: con 3 bits, $[0{,}125;\ 0{,}25)$ empieza antes de $0{,}1445$ y no cabe. Con 4 bits, $J(\bin{0011})=[0{,}1875;\ 0{,}25)$ cabe. Palabra \bin{0011}, 4 bits.
\end{enumerate}
\resultado \textbf{Huffman: 4 bits; aritmética: 4 bits (\bin{0011}); ideal: $3{,}245$ bits. Con mensajes largos, la aritmética se acerca a $H_b(1/4)=0{,}811$ bits por símbolo (\cref{cor:por-simbolo}) y Huffman se queda en 1.}
\comprobacion $3{,}245/4=0{,}811=H_b(1/4)$ (\cref{eje:1-binaria}): «ABBB» tiene exactamente las frecuencias de la fuente. La construcción del \cref{cor:longitud} daría $\lceil 3{,}245\rceil+1=5$ bits.
\end{solucion}

\begin{solucion}{eje:4-bpb}
\datos $\mathcal{J}=0{,}7$ nats por token; $\bar b=3{,}5$ bytes por token.
\planteamiento $\bpb=\mathcal{J}/(\bar b\ln 2)$ (\cref{prop:bpb}); factor de compresión $=8/\bpb$.
\pasos
\begin{enumerate}[label=\arabic*.]
  \item Denominador: $3{,}5\cdot 0{,}6931=2{,}426$.
  \item $\bpb=0{,}7/2{,}426=0{,}2885$ bits por byte.
  \item Factor: $8/0{,}2885=27{,}73$.
\end{enumerate}
\resultado \textbf{$0{,}289$ bits por byte; compresión de unas $27{,}7$ veces.}
\comprobacion por pasos: $0{,}7$ nats $=0{,}7\cdot 1{,}443=1{,}010$ bits por token; dividido entre $3{,}5$ bytes da $0{,}2886$.
\end{solucion}

\begin{solucion}{eje:4-equivocado}
\datos $p=(1/4,1/4,1/2)$; $q=(1/2,1/4,1/4)$; $\ell(x)=-\log q(x)$.
\planteamiento calcular $\E_p[\ell]$, $H(p)$ y $\KL{p}{q}$ por separado y comparar.
\pasos
\begin{enumerate}[label=\arabic*.]
  \item Longitudes: $\ell(\mathrm{a})=1$, $\ell(\mathrm{b})=2$, $\ell(\mathrm{c})=2$.
  \item $\E_p[\ell]=\tfrac14\cdot 1+\tfrac14\cdot 2+\tfrac12\cdot 2=0{,}25+0{,}5+1=1{,}75$ bits.
  \item $H(p)=\tfrac14\cdot 2+\tfrac14\cdot 2+\tfrac12\cdot 1=1{,}5$ bits.
  \item $\KL{p}{q}=\tfrac14\log\tfrac{1/4}{1/2}+\tfrac14\log 1+\tfrac12\log\tfrac{1/2}{1/4}=-0{,}25+0+0{,}5=0{,}25$ bits.
\end{enumerate}
\resultado \textbf{$\E_p[\ell]=1{,}75$ bits $=H(p)+\KL{p}{q}=1{,}5+0{,}25$: el modelo equivocado desperdicia $0{,}25$ bits por símbolo.}
\comprobacion el código óptimo para $p$ tendría longitudes $(2,2,1)$ y longitud media $1{,}5$ bits; el exceso es la divergencia, como en la \cref{prop:descomposicion}.
\end{solucion}

\begin{solucion}{eje:4-equivocado-teo}
\datos $\ell(x)=\lceil-\log q(x)\rceil$; $X\sim p$; $q(x)>0$ si $p(x)>0$.
\planteamiento acotar $\ell(x)$ entre $-\log q(x)$ y $-\log q(x)+1$, tomar esperanzas con $p$ y usar $H(p,q)=H(p)+\KL{p}{q}$.
\pasos
\begin{enumerate}[label=\arabic*.]
  \item Para todo $y$, $y\le\lceil y\rceil<y+1$; con $y=-\log q(x)$: $-\log q(x)\le\ell(x)<-\log q(x)+1$.
  \item Multiplicando por $p(x)\ge 0$ y sumando sobre $x$ (la desigualdad estricta se conserva porque $\sum_xp(x)=1>0$): $H(p,q)\le\E_p[\ell(X)]<H(p,q)+1$.
  \item Por la \cref{prop:descomposicion}, $H(p,q)=H(p)+\KL{p}{q}$.
  \item Las longitudes cumplen Kraft (véase tras la \cref{def:shannon}, con $q$ en lugar de $p$), luego el código existe.
\end{enumerate}
\resultado \textbf{$H(p)+\KL{p}{q}\le\E_p[\ell(X)]<H(p)+\KL{p}{q}+1$: usar el modelo $q$ cuesta $\KL{p}{q}$ bits por símbolo adicionales.}
\comprobacion con los datos del \cref{eje:4-equivocado}, las longitudes de Shannon son exactas y $\E_p[\ell]=1{,}75$ coincide con la cota inferior $1{,}5+0{,}25$.
\end{solucion}

% ---------------------------------------------------------------------
\section{Tema 5}

\begin{solucion}{eje:5-beta}
\datos $\pi^{*}(y)=\pref(y)e^{r(y)/\beta}/Z$; $y^{\star}$ es la única respuesta con recompensa máxima $r_{\max}$.
\planteamiento dividir numerador y denominador por $e^{r_{\max}/\beta}$ y estudiar los exponentes.
\pasos
\begin{enumerate}[label=\arabic*.]
  \item $\pi^{*}(y)=\dfrac{\pref(y)\,e^{(r(y)-r_{\max})/\beta}}{\sum_{y'}\pref(y')\,e^{(r(y')-r_{\max})/\beta}}$.
  \item $\beta\to\infty$: todos los exponentes tienden a~0 y cada exponencial a~1, de modo que $\pi^{*}(y)\to\pref(y)/\sum_{y'}\pref(y')=\pref(y)$.
  \item $\beta\to 0^{+}$: para $y\ne y^{\star}$, $r(y)-r_{\max}<0$ y $e^{(r(y)-r_{\max})/\beta}\to 0$; para $y^{\star}$ el exponente es~0. Luego $\pi^{*}(y^{\star})\to\pref(y^{\star})/\pref(y^{\star})=1$.
  \item En el ajuste de un LLM, $\beta$ regula el compromiso: grande, el modelo apenas se aparta del modelo supervisado; pequeño, persigue la recompensa y queda expuesto a explotar los defectos del modelo de recompensa.
\end{enumerate}
\resultado \textbf{$\beta\to\infty$: $\pi^{*}\to\pref$; $\beta\to 0^{+}$: $\pi^{*}$ se concentra en la respuesta de recompensa máxima. $\beta$ actúa como una temperatura.}
\comprobacion es la tendencia de la tabla del \cref{ej:rlhf}: la divergencia pasa de $0{,}915$ a $0{,}084$ nats al subir $\beta$ de $0{,}5$ a~2.
\end{solucion}

\begin{solucion}{eje:5-monedas}
\datos $n=10$; macroestados $k=5$ y $k=2$ caras.
\planteamiento $W=\binom{n}{k}$ microestados equiprobables; $S/k_B$ en bits es $\log W$; la cota es $nH_b(k/n)$.
\pasos
\begin{enumerate}[label=\arabic*.]
  \item $\binom{10}{5}=252$ y $\log 252=\ln 252/\ln 2=5{,}529/0{,}6931=7{,}977$ bits.
  \item Cota: $10\,H_b(0{,}5)=10$ bits.
  \item $\binom{10}{2}=45$ y $\log 45=3{,}807/0{,}6931=5{,}492$ bits.
  \item Cota: $10\,H_b(0{,}2)=10\cdot 0{,}7219=7{,}219$ bits.
\end{enumerate}
\resultado \textbf{«5 caras»: $7{,}977$ bits frente a la cota de 10; «2 caras»: $5{,}492$ bits frente a $7{,}219$.}
\comprobacion ambas cumplen la \cref{prop:binomial}. Con $n=10$ la cota aún dista de 2 bits; con $n=100$ y $k=50$ la diferencia es $100-96{,}35=3{,}65$ bits, mucho menor en proporción: el exponente $nH_b$ es el término dominante.
\end{solucion}

\begin{solucion}{eje:5-landauer}
\datos $8\cdot 10^{9}$ bits; $T=300$~K; $k_B=1{,}380649\cdot10^{-23}$~J/K; $1$~eV $=1{,}602\cdot10^{-19}$~J.
\planteamiento calor mínimo por bit $k_BT\ln 2$ (\cref{prop:landauer}), multiplicado por el número de bits.
\pasos
\begin{enumerate}[label=\arabic*.]
  \item Por bit: $1{,}380649\cdot10^{-23}\cdot 300\cdot 0{,}6931=2{,}871\cdot10^{-21}$~J.
  \item Total: $8\cdot10^{9}\cdot 2{,}871\cdot10^{-21}=2{,}297\cdot10^{-11}$~J.
  \item En electronvoltios: $2{,}871\cdot10^{-21}/1{,}602\cdot10^{-19}=0{,}01792$~eV por bit.
\end{enumerate}
\resultado \textbf{unos $2{,}30\cdot10^{-11}$~J para todo el gigabyte; unos $0{,}0179$~eV por bit.}
\comprobacion $k_BT$ a 300~K es aproximadamente $0{,}0259$~eV, y $0{,}0259\cdot 0{,}6931=0{,}0179$~eV.
\end{solucion}

\begin{solucion}{eje:5-rlhf}
\datos $\pref=(1/2,1/2)$; $r=(0,1)$; $\beta=1$.
\planteamiento $\pi^{*}(y)=\pref(y)e^{r(y)}/Z$ con $Z=\sum_y\pref(y)e^{r(y)}$; el valor óptimo es $\ln Z$ (\cref{teo:variacional}).
\pasos
\begin{enumerate}[label=\arabic*.]
  \item $Z=\tfrac12e^{0}+\tfrac12e^{1}=(1+2{,}718)/2=1{,}859$.
  \item $\pi^{*}_1=0{,}5/1{,}859=0{,}2689$; $\pi^{*}_2=0{,}5\cdot 2{,}718/1{,}859=0{,}7311$.
  \item $\E_{\pi^{*}}[r]=0\cdot 0{,}2689+1\cdot 0{,}7311=0{,}7311$.
  \item $\KL{\pi^{*}}{\pref}=0{,}2689\ln\tfrac{0{,}2689}{0{,}5}+0{,}7311\ln\tfrac{0{,}7311}{0{,}5}=0{,}2689\cdot(-0{,}6203)+0{,}7311\cdot 0{,}3799=-0{,}1668+0{,}2777=0{,}1109$ nats.
  \item $\mathcal{F}(\pi^{*})=0{,}7311-0{,}1109=0{,}6202$.
  \item $\ln Z=\ln 1{,}859=0{,}6201$.
\end{enumerate}
\resultado \textbf{$\pi^{*}\approx(0{,}269;\ 0{,}731)$; $\E[r]\approx 0{,}731$; divergencia $\approx 0{,}111$ nats; $\mathcal{F}(\pi^{*})\approx 0{,}620=\ln Z$.}
\comprobacion $\pi^{*}$ suma~1, y $\pi^{*}_2=e/(1+e)=\sigma(1)$, la sigmoide de la diferencia de recompensas. La discrepancia en la cuarta cifra se debe al redondeo intermedio.
\end{solucion}

\begin{solucion}{eje:5-lagrange}
\datos $\X$ finito; $f\colon\X\to\R$; restricciones $\sum_xp(x)=1$ y $\sum_xp(x)f(x)=\alpha$; $p(x)>0$ para todo $x$.
\planteamiento maximizar $H(p)=-\sum_xp(x)\ln p(x)$ con dos restricciones de igualdad mediante la lagrangiana (módulo~4).
\pasos
\begin{enumerate}[label=\arabic*.]
  \item Lagrangiana: $\mathcal{L}(p,\mu,\lambda)=-\sum_xp(x)\ln p(x)+\mu\Bigl(\sum_xp(x)-1\Bigr)-\lambda\Bigl(\sum_xp(x)f(x)-\alpha\Bigr)$.
  \item Condición de estacionariedad para cada $x$: $\dfrac{\partial\mathcal{L}}{\partial p(x)}=-\ln p(x)-1+\mu-\lambda f(x)=0$.
  \item Despejando, $p(x)=e^{\mu-1}e^{-\lambda f(x)}$.
  \item La normalización da $e^{\mu-1}\sum_xe^{-\lambda f(x)}=1$, es decir, $e^{\mu-1}=1/Z$ con $Z=\sum_xe^{-\lambda f(x)}$. Así, $p(x)=e^{-\lambda f(x)}/Z$; el valor de $\lambda$ queda fijado por la restricción $\sum_xp(x)f(x)=\alpha$.
  \item Comparación: el método de Lagrange solo da una condición necesaria (punto crítico, y con la hipótesis $p>0$ para poder derivar). Para concluir que es un máximo global hay que añadir que $H$ es estrictamente cóncava y las restricciones lineales. La demostración del \cref{teo:maxent} con la desigualdad de Gibbs da directamente la optimalidad global y la unicidad, y muestra cuánto pierde cualquier otra $p$: exactamente $\KL{p}{p_\lambda}$.
\end{enumerate}
\resultado \textbf{la distribución de máxima entropía con media de $f$ fijada es $p(x)=e^{-\lambda f(x)}/Z$, una distribución de Boltzmann.}
\comprobacion con $f$ constante la restricción no informa y debe salir la uniforme: en efecto, $e^{-\lambda f}$ es constante y $p(x)=1/\abs{\X}$, de entropía máxima $\ln\abs{\X}$.
\end{solucion}

\begin{solucion}{eje:5-dpo}
\datos $r(x,y)=\beta\ln\frac{\pi^{*}(y\mid x)}{\pref(y\mid x)}+\beta\ln Z(x)$; $P(y_1\succ y_2\mid x)=\sigma\bigl(r(x,y_1)-r(x,y_2)\bigr)$; en (b), $\pref=(1/2,1/2)$, $\pi^{*}=(0{,}2689;\ 0{,}7311)$ y $\beta=1$.
\planteamiento restar las recompensas de las dos respuestas para el mismo mensaje $x$.
\pasos
\begin{enumerate}[label=\arabic*.]
  \item $r(x,y_1)-r(x,y_2)=\beta\ln\frac{\pi^{*}(y_1\mid x)}{\pref(y_1\mid x)}+\beta\ln Z(x)-\beta\ln\frac{\pi^{*}(y_2\mid x)}{\pref(y_2\mid x)}-\beta\ln Z(x)$.
  \item Los términos $\beta\ln Z(x)$ se cancelan, porque ambas respuestas responden al mismo mensaje:
  $P(y_1\succ y_2\mid x)=\sigma\Bigl(\beta\ln\frac{\pi^{*}(y_1\mid x)}{\pref(y_1\mid x)}-\beta\ln\frac{\pi^{*}(y_2\mid x)}{\pref(y_2\mid x)}\Bigr)$.
  \item Por tanto, la verosimilitud de las comparaciones humanas puede escribirse en función de la política, sin calcular $Z(x)$ ni entrenar un modelo de recompensa: es la base de DPO.
  \item (b) $r_2-r_1=\ln\frac{0{,}7311}{0{,}5}-\ln\frac{0{,}2689}{0{,}5}=0{,}3799-(-0{,}6203)=1{,}000$.
  \item $P(y_2\succ y_1)=\sigma(1)=1/(1+e^{-1})=0{,}7311$.
\end{enumerate}
\resultado \textbf{(a) la probabilidad de preferencia solo depende de $\beta\ln(\pi^{*}/\pref)$ de cada respuesta; (b) diferencia implícita de recompensas igual a~1 y $P(y_2\succ y_1)\approx 0{,}731$.}
\comprobacion la diferencia recuperada, $1$, es la diferencia de recompensas $r=(0,1)$ del \cref{eje:5-rlhf}: la política óptima codifica la recompensa salvo una constante por mensaje.
\end{solucion}
